% Options for packages loaded elsewhere
\PassOptionsToPackage{unicode}{hyperref}
\PassOptionsToPackage{hyphens}{url}
\documentclass[
]{article}
\usepackage{xcolor}
\usepackage[margin=1in]{geometry}
\usepackage{amsmath,amssymb}
\setcounter{secnumdepth}{5}
\usepackage{iftex}
\ifPDFTeX
  \usepackage[T1]{fontenc}
  \usepackage[utf8]{inputenc}
  \usepackage{textcomp} % provide euro and other symbols
\else % if luatex or xetex
  \usepackage{unicode-math} % this also loads fontspec
  \defaultfontfeatures{Scale=MatchLowercase}
  \defaultfontfeatures[\rmfamily]{Ligatures=TeX,Scale=1}
\fi
\usepackage{lmodern}
\ifPDFTeX\else
  % xetex/luatex font selection
\fi
% Use upquote if available, for straight quotes in verbatim environments
\IfFileExists{upquote.sty}{\usepackage{upquote}}{}
\IfFileExists{microtype.sty}{% use microtype if available
  \usepackage[]{microtype}
  \UseMicrotypeSet[protrusion]{basicmath} % disable protrusion for tt fonts
}{}
\makeatletter
\@ifundefined{KOMAClassName}{% if non-KOMA class
  \IfFileExists{parskip.sty}{%
    \usepackage{parskip}
  }{% else
    \setlength{\parindent}{0pt}
    \setlength{\parskip}{6pt plus 2pt minus 1pt}}
}{% if KOMA class
  \KOMAoptions{parskip=half}}
\makeatother
\usepackage{color}
\usepackage{fancyvrb}
\newcommand{\VerbBar}{|}
\newcommand{\VERB}{\Verb[commandchars=\\\{\}]}
\DefineVerbatimEnvironment{Highlighting}{Verbatim}{commandchars=\\\{\}}
% Add ',fontsize=\small' for more characters per line
\usepackage{framed}
\definecolor{shadecolor}{RGB}{248,248,248}
\newenvironment{Shaded}{\begin{snugshade}}{\end{snugshade}}
\newcommand{\AlertTok}[1]{\textcolor[rgb]{0.94,0.16,0.16}{#1}}
\newcommand{\AnnotationTok}[1]{\textcolor[rgb]{0.56,0.35,0.01}{\textbf{\textit{#1}}}}
\newcommand{\AttributeTok}[1]{\textcolor[rgb]{0.13,0.29,0.53}{#1}}
\newcommand{\BaseNTok}[1]{\textcolor[rgb]{0.00,0.00,0.81}{#1}}
\newcommand{\BuiltInTok}[1]{#1}
\newcommand{\CharTok}[1]{\textcolor[rgb]{0.31,0.60,0.02}{#1}}
\newcommand{\CommentTok}[1]{\textcolor[rgb]{0.56,0.35,0.01}{\textit{#1}}}
\newcommand{\CommentVarTok}[1]{\textcolor[rgb]{0.56,0.35,0.01}{\textbf{\textit{#1}}}}
\newcommand{\ConstantTok}[1]{\textcolor[rgb]{0.56,0.35,0.01}{#1}}
\newcommand{\ControlFlowTok}[1]{\textcolor[rgb]{0.13,0.29,0.53}{\textbf{#1}}}
\newcommand{\DataTypeTok}[1]{\textcolor[rgb]{0.13,0.29,0.53}{#1}}
\newcommand{\DecValTok}[1]{\textcolor[rgb]{0.00,0.00,0.81}{#1}}
\newcommand{\DocumentationTok}[1]{\textcolor[rgb]{0.56,0.35,0.01}{\textbf{\textit{#1}}}}
\newcommand{\ErrorTok}[1]{\textcolor[rgb]{0.64,0.00,0.00}{\textbf{#1}}}
\newcommand{\ExtensionTok}[1]{#1}
\newcommand{\FloatTok}[1]{\textcolor[rgb]{0.00,0.00,0.81}{#1}}
\newcommand{\FunctionTok}[1]{\textcolor[rgb]{0.13,0.29,0.53}{\textbf{#1}}}
\newcommand{\ImportTok}[1]{#1}
\newcommand{\InformationTok}[1]{\textcolor[rgb]{0.56,0.35,0.01}{\textbf{\textit{#1}}}}
\newcommand{\KeywordTok}[1]{\textcolor[rgb]{0.13,0.29,0.53}{\textbf{#1}}}
\newcommand{\NormalTok}[1]{#1}
\newcommand{\OperatorTok}[1]{\textcolor[rgb]{0.81,0.36,0.00}{\textbf{#1}}}
\newcommand{\OtherTok}[1]{\textcolor[rgb]{0.56,0.35,0.01}{#1}}
\newcommand{\PreprocessorTok}[1]{\textcolor[rgb]{0.56,0.35,0.01}{\textit{#1}}}
\newcommand{\RegionMarkerTok}[1]{#1}
\newcommand{\SpecialCharTok}[1]{\textcolor[rgb]{0.81,0.36,0.00}{\textbf{#1}}}
\newcommand{\SpecialStringTok}[1]{\textcolor[rgb]{0.31,0.60,0.02}{#1}}
\newcommand{\StringTok}[1]{\textcolor[rgb]{0.31,0.60,0.02}{#1}}
\newcommand{\VariableTok}[1]{\textcolor[rgb]{0.00,0.00,0.00}{#1}}
\newcommand{\VerbatimStringTok}[1]{\textcolor[rgb]{0.31,0.60,0.02}{#1}}
\newcommand{\WarningTok}[1]{\textcolor[rgb]{0.56,0.35,0.01}{\textbf{\textit{#1}}}}
\usepackage{longtable,booktabs,array}
\newcounter{none} % for unnumbered tables
\usepackage{calc} % for calculating minipage widths
% Correct order of tables after \paragraph or \subparagraph
\usepackage{etoolbox}
\makeatletter
\patchcmd\longtable{\par}{\if@noskipsec\mbox{}\fi\par}{}{}
\makeatother
% Allow footnotes in longtable head/foot
\IfFileExists{footnotehyper.sty}{\usepackage{footnotehyper}}{\usepackage{footnote}}
\makesavenoteenv{longtable}
\usepackage{graphicx}
\makeatletter
\newsavebox\pandoc@box
\newcommand*\pandocbounded[1]{% scales image to fit in text height/width
  \sbox\pandoc@box{#1}%
  \Gscale@div\@tempa{\textheight}{\dimexpr\ht\pandoc@box+\dp\pandoc@box\relax}%
  \Gscale@div\@tempb{\linewidth}{\wd\pandoc@box}%
  \ifdim\@tempb\p@<\@tempa\p@\let\@tempa\@tempb\fi% select the smaller of both
  \ifdim\@tempa\p@<\p@\scalebox{\@tempa}{\usebox\pandoc@box}%
  \else\usebox{\pandoc@box}%
  \fi%
}
% Set default figure placement to htbp
\def\fps@figure{htbp}
\makeatother
\setlength{\emergencystretch}{3em} % prevent overfull lines
\providecommand{\tightlist}{%
  \setlength{\itemsep}{0pt}\setlength{\parskip}{0pt}}
\usepackage{bookmark}
\IfFileExists{xurl.sty}{\usepackage{xurl}}{} % add URL line breaks if available
\urlstyle{same}
\hypersetup{
  pdftitle={Capr Design Guide and Roadmap},
  hidelinks,
  pdfcreator={LaTeX via pandoc}}

\title{Capr Design Guide and Roadmap}
\author{}
\date{\vspace{-2.5em}}

\begin{document}
\maketitle

{
\setcounter{tocdepth}{2}
\tableofcontents
}
\begin{Shaded}
\begin{Highlighting}[]
\FunctionTok{library}\NormalTok{(Capr)}
\end{Highlighting}
\end{Shaded}

\section{Executive Summary}\label{executive-summary}

Capr, pronounced `kay-pr' like the edible flower, is an R package that
serves as a programmatic interface to defining cohort definitions for
data standardized to the OMOP CDM. Capr allows for templating cohort
logic, where a user can shuffle through concept sets to initiate
variations of a cohort definition. Cohort templating is a critical value
add that Capr provides to OHDSI software. The intent of this software is
to make the process of defining cohorts less of a clerical task (via
point and click) and allow for OHDSI studies to be designed in a single
development environment. The value of developing Capr for Odysseus is
that it integrates cohort development into the development pipeline,
enables templating which can be leveraged in a phenotyping framework.

The mechanics of Capr are based on emulating the representation of a
cohort definition as specified by circe-be. Circe-be is the underlying
mechanism that builds sql based on input specification defined by a json
structure that fills in a java class. Capr can be used now be used
without a connection to a CDM, however the user must have knowledge of
the OMOP CDM providing concept ids as inputs. Capr returns an R object
that can be converted to a json string. Commonly, cohort definitions are
distributed within the OHDSI community as json files, therefore it is a
natural endpoint for Capr, allowing circe to serialize to ohdsiSql. The
functionality of Capr includes: creation of a concept set, specifying
the cohort logic, building a cohort definition and coercing it to a json
representation. Future functionality will provide print and validity
methods for class, importing Capr from json and saving/loading cohort
sub-components to disk. Capr also has room to grow in terms of providing
functionality to describe a cohort definition both by graph and text,
modifying cohort definitions by addition of Capr objects and enumerating
sub-components.

Currently, Capr is preparing for a release of v2.0.0 to the OHDSI
community github as a HADES R package. The release date is anticipated
for early spring 2023. There is a plan to incorporate complimentary
changes to v2.0.0 by late summer 2023 in addition to suggestions based
on user feedback.

\section{Introduction}\label{introduction}

\subsection{Purpose}\label{purpose}

Capr is a programmatic interface to defining cohort definition for the
OMOP CDM. Capr builds a representation of a cohort through R that can be
serialized to json. The json structure can be imported into ATLAS. This
structure follows that of circe-be, that object representation of cohort
definitions in OHDSI. A valid circe-be json can be serialized into
ohdsisql, a parameterized version of sql that builds cohort under the
assumptions of the OMOP CDM structure. Capr fills a gap in OHDSI
software in providing a programmatic tool for cohort construction as an
alternative to ATLAS. A programmatic tool enables us to parameterize
definitions. This has a few advantages:

\begin{enumerate}
\def\labelenumi{\arabic{enumi}.}
\tightlist
\item
  \emph{reduce clerical error} -- minimize copy/paste activities in
  ATLAS by instead offering a cohort skeleton and iterating across
  concept sets.
\item
  \emph{Focus on concept set development} -- our ability to properly
  enumerate a cohort of patients ultimately relies on whether the
  correct concepts have been identified for a database. Many cohort
  definition skeletons are simple where the variation in the definition
  comes largely from the concepts. Capr allows us to expediate the
  definition process and focus on identifying the rights concepts.
\item
  \emph{Provide alternative entry points to OHDSI analysis} -- ATLAS
  installation is not guaranteed in a study site with OMOP data. Further
  some researchers do not like to use a GUI like ATLAS and prefer to use
  a programming language like R. Capr builds an equivalent cohort
  definition to ATLAS meaning there is an alternative entry point to an
  OHDSI study. Capr allows a researcher to remain in a single interface
  to complete their entire analysis: from cohort definition to
  diagnostics to analysis. It is however, suggested the Capr be used in
  tandem with ATLAS for clinical review.
\end{enumerate}

\subsection{User-base}\label{user-base}

A Capr user is someone who has knowledge of OHDSI and the OMOP CDM and
is comfortable programming in R. Capr is intended to be used by data
scientists to help create a scripted representation of a cohort
definition that has been defined and reviewed by a clinical team.
Defining cohorts is a team effort, however Capr gives programmatically
inclined users a tool to produce cohort definitions for their clinical
team.

\section{Design}\label{design}

\subsection{A. Orientation}\label{a.-orientation}

\subsubsection{i. Entry Point}\label{i.-entry-point}

Ideally, Capr should be used in tandem with a connection to an OMOP
vocabulary. Connection to an OMOP vocabulary can be done through
DatabaseConnector or CDMConnector. This allows users to lookup concepts
before applying them in Capr objects. Capr does not provide
functionality to optimize concept lookup; this is a feature passed to
other OHDSI software. Capr assumes the list of concepts as inputs.

\subsubsection{ii. Exit Point}\label{ii.-exit-point}

Capr stops upon development of the json structure, which is typically
saved as a file in a directory. The previous version of Capr allowed for
the user to create the ohdsiSql via CirceR. However, this was deemed an
overlap between other OHDSI R packages (CirceR, CohortGenerator etc.)
and it felt more appropriate for the Capr process to end at building the
json file. Ultimately the item that is passed across the OHDSI network
is the cohort definition json. The json file can be upload into ATLAS or
into R where the ohdsiSql can be built using CirceR. Json is a straight
input to build the SQL, so it was deemed the minimally sufficient output
needed. We intend to draw distinct lines between OHDSI software in order
to improve standard workflows, that are either site internal or for
network studies.

\subsection{B. Functionality}\label{b.-functionality}

Capr provides the programmatic interface for constructing cohort
definitions that maps to the circe-be logic. In order to maintain this
correspondence, the functionality of Capr covers: building concept sets,
composing and building the cohort definition, and coercing Capr to json.
Future functionality will be added into the section once they are
released.

\subsubsection{i. Concept Set Builder}\label{i.-concept-set-builder}

Before defining cohort logic, the user must specify concept sets. Capr
provides functionality to build concept sets off standardized OMOP
concept IDs. How the user finds these IDs is left to other tools. A
concept set is built in Capr as follows:

\begin{Shaded}
\begin{Highlighting}[]
\FunctionTok{cs}\NormalTok{(}\FunctionTok{descendants}\NormalTok{(}\DecValTok{201826}\DataTypeTok{L}\NormalTok{), }\AttributeTok{name =} \StringTok{"T2D"}\NormalTok{)}
\FunctionTok{cs}\NormalTok{(}\DecValTok{201826}\DataTypeTok{L}\NormalTok{, }\DecValTok{201254}\DataTypeTok{L}\NormalTok{, }\AttributeTok{name =} \StringTok{"Diabetes"}\NormalTok{)}
\end{Highlighting}
\end{Shaded}

The first example in code snippet 1 provides the concept set and
specifies that it includes all descendants. Capr supplies a command for
the three types of concept mapping: include descendants, exclude, and
include mapped. The second example in code snippet 1 provides builds a
concept set based on a vector of concept IDs. The first argument of the
cs command is an ellipse capturing a variety of input giving the user
flexibility. The second argument gives the concept set a name, which in
general should be supplied for readability.

When a concept set is created in Capr, it is tagged with a global unique
identifier (GUID) using the R package digest. The purpose of this is to
ensure that the id of concept set is unique to its construction. In
ATLAS, the concept id attached to the concept set is a local id. This
means that the id is an integer starting from 0 that identifies the
concept set in the definition. This creates an issue if we were to
parameterize the concept set, therefore the tag must be unique to the
concept set. Since circe-be uses integers, Capr must resolve this issue
on the backend, something described in the coercion section.

\subsubsection{ii. Cohort Builder}\label{ii.-cohort-builder}

Next, a user must be able to build a cohort definition in R, therefore
Capr provides functionality to define a cohort definition and
sub-components of the cohort definition. Capr provides a constructor for
each circe-be element (i.e.~query, count, group, windows). A cohort is
defined using the function:

\begin{Shaded}
\begin{Highlighting}[]
\FunctionTok{cohort}\NormalTok{(entry, attrition, exit, era)}
\end{Highlighting}
\end{Shaded}

The inputs for the cohort function are the four components for defining
a cohort: entry which accrues the initial set of patients, attrition
which defines the logic that filters the initial set of patients based
on inclusion and exclusion criteria, exit which determines how the
person exits the cohort and era which defines the span of time in which
events are considered to have a cohort event. Each of these four
elements are constructed based on circe-be subcomponents including:

\begin{enumerate}
\def\labelenumi{\arabic{enumi})}
\tightlist
\item
  \emph{Query}: an object that defines which concepts to extract from a
  domain table in the CDM;
\item
  \emph{Attribute}: an object that specifies a subpopulation of the
  query;
\item
  \emph{Criteria}: an object that temporally enumerates the presence or
  absence of an event relative to a point in time; and
\item
  \emph{Group}: an object that binds multiple criteria and sub-groups as
  a single expression.
\end{enumerate}

\paragraph{1) Query}\label{query}

A query is constructed via the name of the domain table in the OMOP CDM.
For example:

\begin{Shaded}
\begin{Highlighting}[]
\FunctionTok{visit}\NormalTok{(conceptSet, ...)}
\FunctionTok{conditionOccurrence}\NormalTok{(conceptSet, ...)}
\end{Highlighting}
\end{Shaded}

The conceptSet input takes on a a concept set built from the cs command.
The dots argument is where the user can supply attributes that
contextualize the query to a specific subpopulation.

\paragraph{2) Attribute}\label{attribute}

An attribute has its own set of constructor functions that are
variations on four types of attributes that may be used:

\begin{itemize}
\tightlist
\item
  \textbf{op}: a boolean operator matched with either a numeric, integer
  or date value
\item
  \textbf{concept}: attributed built on an OMOP concept id
\item
  \textbf{logic} an attribute defined by a \texttt{T/F} toggle
\item
  \textbf{nested}: an attribute based on a Capr group object
\end{itemize}

An example of using an attribute within a query is shown in code snippet
4:

\begin{Shaded}
\begin{Highlighting}[]
\FunctionTok{conditionOccurrence}\NormalTok{(}\FunctionTok{cs}\NormalTok{(}\FunctionTok{descendants}\NormalTok{(}\DecValTok{201826}\DataTypeTok{L}\NormalTok{), }\AttributeTok{name =} \StringTok{"T2D"}\NormalTok{), }\FunctionTok{male}\NormalTok{())}
\FunctionTok{conditionOccurrence}\NormalTok{(}\FunctionTok{cs}\NormalTok{(}\FunctionTok{descendants}\NormalTok{(}\DecValTok{201826}\DataTypeTok{L}\NormalTok{), }\AttributeTok{name =} \StringTok{"T2D"}\NormalTok{), }\FunctionTok{age}\NormalTok{(}\FunctionTok{gte}\NormalTok{(}\DecValTok{18}\NormalTok{)))}
\FunctionTok{conditionOccurrence}\NormalTok{(}\FunctionTok{cs}\NormalTok{(}\FunctionTok{descendants}\NormalTok{(}\DecValTok{201826}\DataTypeTok{L}\NormalTok{), }\AttributeTok{name =} \StringTok{"T2D"}\NormalTok{), }\FunctionTok{firstOccurrence}\NormalTok{())}
\end{Highlighting}
\end{Shaded}

The first example in code snippet 4 utilizes a concept attribute. The
\texttt{male()} function automatically maps the concept to the male
concept id. Other concept attribute functions will have mappers and
helpers to assist the user in selecting the correct concept to add as an
attribute. The second example of an attribute in code snippet 4 is for
an op attribute. Any op attribute begins with the domain (i.e.~age,
startDate etc) and then takes an input of an op. An op is a function
with the following signatures:

{\def\LTcaptype{none} % do not increment counter
\begin{longtable}[]{@{}ll@{}}
\toprule\noalign{}
Operator & Name \\
\midrule\noalign{}
\endhead
\bottomrule\noalign{}
\endlastfoot
gt & Greater than \\
gte & Greater than or equal to \\
lt & Less than \\
lte & Less than or equal to \\
eq & Equal to \\
bt & Between \\
nbt & Not between \\
\end{longtable}
}

Note that with the bt and nbt functions the input requires two values, a
start and stop, specifying a bounded range of values. The last example
in code snippet 4 is for a logic attribute. The logic attribute is
simple in that if it is added to the query it will toggle TRUE for a the
specification of the domain. In the example we omit one for a nested
attribute because it is more complex and requires explanation of a
criteria and group, which are stated in the upcoming sections. A full
list of available attributes will be included in a future appendix.

\paragraph{3) Criteria}\label{criteria}

A criteria is constructed first on an operator that specifies the number
of times the criteria is observed, second the query of which the
criteria is based on and finally an aperture (i.e.~a time window)
specifying the timing relative to the index event at which we are to
observe the specified criteria. Code snippet 5 shows an example of a
criteria call:

\begin{Shaded}
\begin{Highlighting}[]
\FunctionTok{exactly}\NormalTok{(}\DecValTok{0}\NormalTok{, }
  \AttributeTok{query =} \FunctionTok{conditionOccurrence}\NormalTok{(}\FunctionTok{cs}\NormalTok{(}\FunctionTok{descendants}\NormalTok{(}\DecValTok{201826}\DataTypeTok{L}\NormalTok{), }\AttributeTok{name =} \StringTok{"T2D"}\NormalTok{)),}
        \AttributeTok{aperture =} \FunctionTok{duringInterval}\NormalTok{(}\FunctionTok{eventStarts}\NormalTok{(}\SpecialCharTok{{-}}\ConstantTok{Inf}\NormalTok{, }\SpecialCharTok{{-}}\DecValTok{1}\NormalTok{))}
\NormalTok{        )}
\end{Highlighting}
\end{Shaded}

A note for the criteria is that its meaning is tethered to an event
specified in another part of a cohort definition. This means that a
criteria object applies either for cohort attrition (inclusion rules),
additional criteria for the primary criteria or in a nested attribute
where the index event is the query to which the attribute is attached.
The aperture is what restricts the criteria to certain situations. An
aperture is built using the \texttt{duringInterval} command, where the
first option is the start of the window and second is the end of a
window. The temporal deployment of the query makes this object complex.

\paragraph{4) Group}\label{group}

The final sub-component binds multiple criteria and sub-groups in a
single expression. In Capr, this is done using the term with, to which
we attach the options all, any, at most or at least. This specifies for
the candidate to enter the cohort they have to satisfy the
specifications of the group, whether that is at least 1 criteria or all.
In the example below we assume that the query and aperture have been
specified elsewhere in a script.

\begin{Shaded}
\begin{Highlighting}[]
\FunctionTok{withAny}\NormalTok{(}
  \FunctionTok{atLeast}\NormalTok{(}\DecValTok{1}\NormalTok{, }\AttributeTok{query =}\NormalTok{ abLabHb, }\AttributeTok{aperture =}\NormalTok{ ap1),}
  \FunctionTok{atLeast}\NormalTok{(}\DecValTok{1}\NormalTok{, }\AttributeTok{query =}\NormalTok{ abLabRan, }\AttributeTok{aperture =}\NormalTok{ ap1),}
  \FunctionTok{atLeast}\NormalTok{(}\DecValTok{1}\NormalTok{, }\AttributeTok{query =}\NormalTok{ abLabFast, }\AttributeTok{aperture =}\NormalTok{ ap1)}
\NormalTok{)}
\end{Highlighting}
\end{Shaded}

In code snippet 6, we start the group by specifying what needs to be
satisfied in order for the candidate to enter the cohort, in this case
they must satisfy all the queries. If the group is based on
\texttt{atLeast} or \texttt{atMost} then we must also specify an integer
stating how many components need to be satisfied. Groups make complexity
possible in a cohort definition where there are multiple factors that
specify qualification of remaining in a cohort. Groups can also take on
subgroups, for example we could say \texttt{withAll(withAny(..))}. The
group is also the basis of the nested attribute. Specification that a
group is for a nested attribute is as simple as adding
\texttt{nestedWith} to the function signature.

\subsubsection{iii. Coercion/Mapping}\label{iii.-coercionmapping}

For OHDSI studies, the key representation of a cohort definition from a
software perspective is via a json file. Thus, we need to coerce the
Capr object into a json file so that it can be used within studies or
exchanged with other researchers. Previously we had constructed a cohort
definition via an R representation supported by Capr. Now we must turn
this R object into a json file to save to disk. Coercion is a complex
task because Capr must track two aspects of the definition: 1) the
concept sets indicating the codes to search within the CDM and 2) the
logic specifying how to deploy the codes on the CDM. The following
explains how this process is achieved in an abstract sense.

\textbf{Step 1: Collect and replace concept identifier}

Recall that a concept set is initialized with a GUID from the digest R
package. While this identifier is useful in Capr, in circe the id must
be an integer starting from 0. This ID is an internal reference for the
cohort definition. When coercing the Capr structure to a json we need to
collect all the GUIDs in the definition, remove duplicates and enumerate
each unique concept set for the internal reference. Capr does this in
the background using the internal functions \texttt{collectGuid} and
\texttt{replaceCodeSetId}. The first function builds the unique list and
the second goes into each Capr object and replace the id slot with the
new integer. After this process the object remains a Capr S4 cohort
structure but the ids within each query have changed from GUID
(character) to integer.

\textbf{Step 2: Coerce concept set from S4 Capr to S4 list}

Once the concept sets have been set within the query object, Capr will
then go through the entire definition and extract just the concept set
objects and turn them into an S3 list. The internal function
\texttt{listConceptSets} completes this task, returning a list of all
the concept sets used in the definition and removing duplicates.

\textbf{Step 3: Coerce cohort logic from S4 Capr to S3 list}

The cohort logic is now coerced separately from the concept sets. In the
circe-be json the concept sets are a different slot the remainder of the
definition. Capr handles this the same using the internal
\texttt{as.list} to coerce the Capr S4 object to an S3 list. It binds
the concept set list with the cohort logic list.

\textbf{Step 4: Convert S3 list to json}

Once the object is coerced into an S3 list it can be converted into json
using the \texttt{jsonlite} R package. The function
\texttt{jsonlite::toJSON} only works on an S3 list. Capr is designed in
S4 because of the complexity of the system, making it easier to prevent
mistakes by the user.

\subsection{Object Representation}\label{object-representation}

Capr uses S4 classes in order to define the underpinnings of the cohort
definition. These S4 classes are closely related to the circe-be java
classes that specify the structure of the cohort definition. The Capr
classes assume a few differences in order to allow for parameterization
of the cohort definition via templating and provided suggested
organizational improvements to the structure. The decision to use S4
classes was made because of the need to have a strict interface to
adhere to the circe-be class. S4 classes are well suited for very
complex codebases, according to Advanced R. Capr's relation to circe-be
has the sort of complexity, particularly for the coercion aspect of Capr
where the same method is applied slightly differently to each class.
Classes are an important aspect of Capr and S4 provides a good balance
between functional and object-oriented programming.

\subsubsection{Inspecting a cohort
definition}\label{inspecting-a-cohort-definition}

Capr provides \texttt{show()} and \texttt{summary()} methods so a cohort
can be inspected in the console without dumping the raw S4 slots. Typing
an object (or calling \texttt{show()}) gives a compact, one-block
overview; \texttt{summary()} expands the structure, including nested
correlated criteria.

\begin{Shaded}
\begin{Highlighting}[]
\NormalTok{t2dm }\OtherTok{\textless{}{-}} \FunctionTok{cs}\NormalTok{(}\FunctionTok{descendants}\NormalTok{(}\DecValTok{201826}\DataTypeTok{L}\NormalTok{), }\AttributeTok{name =} \StringTok{"T2DM"}\NormalTok{)}
\NormalTok{hf   }\OtherTok{\textless{}{-}} \FunctionTok{cs}\NormalTok{(}\FunctionTok{descendants}\NormalTok{(}\DecValTok{316139}\DataTypeTok{L}\NormalTok{), }\AttributeTok{name =} \StringTok{"Heart failure"}\NormalTok{)}

\NormalTok{cd }\OtherTok{\textless{}{-}} \FunctionTok{cohort}\NormalTok{(}
  \AttributeTok{entry =} \FunctionTok{entry}\NormalTok{(}
    \FunctionTok{conditionOccurrence}\NormalTok{(t2dm, }\FunctionTok{male}\NormalTok{()),}
    \FunctionTok{conditionOccurrence}\NormalTok{(hf, }\FunctionTok{nestedWithAll}\NormalTok{(}\FunctionTok{atLeast}\NormalTok{(}\DecValTok{1}\NormalTok{, }\FunctionTok{conditionOccurrence}\NormalTok{(hf))))}
\NormalTok{  ),}
  \AttributeTok{exit =} \FunctionTok{exit}\NormalTok{(}\AttributeTok{endStrategy =} \FunctionTok{fixedExit}\NormalTok{(}\AttributeTok{offsetDays =} \DecValTok{30}\DataTypeTok{L}\NormalTok{))}
\NormalTok{)}

\NormalTok{cd                                    }\CommentTok{\# compact overview (calls show())}
\FunctionTok{summary}\NormalTok{(cd)                           }\CommentTok{\# fuller view, expands nested criteria}
\FunctionTok{summary}\NormalTok{(cd, }\AttributeTok{listConceptSets =} \ConstantTok{TRUE}\NormalTok{)   }\CommentTok{\# also lists concept sets (id + logic)}
\end{Highlighting}
\end{Shaded}

Passing \texttt{listConceptSets\ =\ TRUE} appends a concept-set section
to the bottom of the summary showing each set's id and per-concept
\texttt{descendants}/\texttt{excluded}/ \texttt{mapped} logic. If the
cohort has been hydrated with \texttt{getConceptSetDetails()} (see the
\emph{Capr-conceptSets} vignette), that section also shows each
concept's name and domain.

\section{Appendix}\label{appendix}

\end{document}
